\section{Canonicity and restriction of geometric affine homology}

\begin{lemma}[The affine tilde-counit comparison is canonical]
\label{lem:cartier-geometric-affine-counit-homology-canonical}
\leanok
\uses{lem:left-exact-original-homology-composition}
\lean{EnvelopeProperty.Cartier.GeometricHomologyNaturality.affineCounitHomology_canonical}
For an affine scheme $\operatorname{Spec}A$ and a complex $K$ of
quasicoherent modules, the comparison obtained by applying sections
to the original tilde-counit homology isomorphism, preceded by the
original tilde unit, is exactly the original canonical comparison
$H^q(\Gamma K)\to\Gamma(H^qK)$.
\end{lemma}
\begin{proof}
Write $L$ for associated sheaf and $G$ for affine sections, and let
$\nu_F$ denote the canonical homology comparison of a left exact
functor. Its proved composition identity gives
$\nu_{L\circ G}=\nu_L\,;\,L(\nu_G)$, where $a;b$ means first
$a$ and then $b$. Naturality for the original counit
$\epsilon:L\circ G\to\mathrm{Id}$ identifies this composite followed
by $\epsilon$ with the homology map of the counit on the complex.
The identity-functor comparison is the identity by uniqueness from
the cycle projection. Since $L$ preserves homology, its canonical
comparison is the original invertible homology-preservation map.
Cancel this map to obtain
\[
 L(\nu_G)\,;\,\epsilon_{H K}
 =\nu_L^{-1}\,;\,H(\epsilon_K).
\]
The right side is precisely the original tilde-counit homology
isomorphism. Apply the original adjunction bijection; its inverse
sends a map to its associated-sheaf image followed by the counit.
The triangle identity therefore identifies $\nu_G$ with the unit
followed by sections of that original isomorphism, as asserted.
\end{proof}

\begin{theorem}[The frozen geometric affine homology comparison is canonical]
\label{thm:cartier-geometric-affine-homology-canonical}
\leanok
\uses{lem:cartier-geometric-affine-counit-homology-canonical,thm:cartier-geometric-affine-homology-sections,lem:left-exact-original-homology-composition}
\lean{EnvelopeProperty.Cartier.GeometricHomologyNaturality.homologySectionsIso_hom_eq}
For every original affine open $W$ and complex $K$ of quasicoherent
modules, the entire forward map of the original geometric affine
homology isomorphism is exactly the original canonical section
homology comparison.
\end{theorem}
\begin{proof}
Let $j:\operatorname{Spec}\Gamma(W)\to X$ be the original affine
open chart. The original affine chart-section isomorphism is the
restriction along the equality $j(\operatorname{Spec}\Gamma(W))=W$;
both constructions have this same forward map on every section.
Naturality of the tilde unit moves the homology of this chart-section
isomorphism past that unit in the original geometric comparison.
The preceding lemma identifies the intervening affine tilde-counit
comparison with the canonical affine section comparison. Exactness
of original open restriction identifies its homology isomorphism
with its canonical comparison. The proved composition identity now
identifies their composite with the canonical comparison for
chart restriction followed by affine sections. Finally, canonical
comparison naturality for the original chart-section natural
isomorphism gives the canonical comparison for evaluation at $W$.
Every map used is the original map; none is selected from a matching
or naturality premise.
\end{proof}

\begin{theorem}[The entire geometric affine homology restriction square]
\label{thm:cartier-geometric-affine-homology-restriction}
\leanok
\uses{thm:cartier-geometric-affine-homology-canonical,prop:scheme-original-native-section-homology-restriction}
\lean{EnvelopeProperty.Cartier.GeometricHomologyNaturality.homologySectionsIso_restrict}
For every inclusion $U\subset V$ of original affine opens, the
original restriction on homology of section complexes followed by
the smaller original affine homology isomorphism equals the larger
original affine homology isomorphism followed by original restriction
of the same global homology sheaf.
\end{theorem}
\begin{proof}
Replace each entire forward isomorphism by its proved original
canonical comparison. The existing native canonical comparison
restriction theorem then gives exactly the asserted square,
including the canonical homology isomorphism for restriction of
scalars along the original section-ring restriction.
\end{proof}
